\section{Held-out instances}\label{sec:holdout}

The comparisons of \cref{sec:results3d,sec:results2d} were read on the validation instances,
on which the runs were monitored and, in two dimensions, on which the results that motivated
the comparison had been seen. A pre-registered evaluation (\cref{tab:runs}) read eight of
their claims again, with the stored networks unchanged, on instances of the same corpora on
which no recorded run or tool had trained or evaluated a network or computed a label, under
rules fixed before any of them was solved. In three dimensions the corpus has
\numHoldThreeTail{} instances beyond the validation set and the training pool
(\cref{sec:methods:data}); the holdout set is \numHoldThreeN{} of them, without the
\numHoldThreeExcl{} that other recorded sets had drawn again. In two dimensions the corpus
has \numHoldTwoTail{} instances beyond the deepest prefix of its training pool that any
recorded run read; the holdout set is \numHoldTwoN{} of them, without the \numHoldTwoExcl{}
that another recorded tool had rebuilt or whose own files were held outside the training
machine. The pre-registration was stamped on 18 October 2026. The evaluation then labelled
the held-out instances with the sparse direct solver, in \numHoldThreeSolves{} solves in
three dimensions and \numHoldTwoSolves{} in two, and evaluated the stored networks on them;
the supervised transformer with the energy term, whose networks were not stored, was not
evaluated.

% generated from the committed records; do not edit
\begin{table}[tbp]
\centering
\footnotesize
\setlength{\tabcolsep}{4pt}
\caption{The claims of the three- and two-dimensional comparisons read again on held-out instances, as the adjudicator of the holdout evaluation recorded them, beside their values on the validation instances.}
\label{tab:holdout}
\resizebox{\ifdim\width>\textwidth\textwidth\else\width\fi}{!}{%
\begin{tabular}{@{}>{\raggedright\arraybackslash}p{0.07\linewidth}>{\raggedright\arraybackslash}p{0.455\linewidth}>{\raggedright\arraybackslash}p{0.14\linewidth}>{\raggedright\arraybackslash}p{0.14\linewidth}>{\raggedright\arraybackslash}p{0.155\linewidth}@{}}
\toprule
Claim & Statement & Validation & Holdout & Outcome \\
\midrule
T1 (3D) & Label-free transformer better than the displacement-supervised transformer on all five metrics, each seed and the seed mean & 20 of 20 & 20 of 20 & replicated \\
T2 (3D) & Graph network: lower displacement error than the label-free transformer on most instances, each seed; higher von Mises error on at least 95\% of the pairs & 84.4\%--92.6\%; 99.9\% & 86.3\%--95.8\%; 100\% & replicated \\
T3 (3D) & Load cases worse than the zero field, displacement-supervised transformer and graph network, each with a positive energy & 126 and 106 of 3,072 & 340 and 237 of 8,232 & replicated \\
\midrule
H1 (2D) & Label-free below the displacement-supervised transformer in relative energy gap, beyond the guard & $-$71.3\% (guard 12.4\%) & $-$71.3\% (guard 11.5\%) & replicated \\
H2a (2D) & Energy-norm-supervised below the displacement-supervised transformer in relative energy gap, beyond the guard & $-$78.2\% (guard 16.4\%) & $-$42.9\% (guard 47.4\%) & \hspace{0pt}not replicated \\
H2b (2D) & The same in von Mises error & $-$54.9\% (guard 12.3\%) & $-$52.8\% (guard 12.9\%) & replicated \\
H2c (2D) & The same, by more than half & $-$54.9\% & $-$52.8\% & replicated \\
D2 (2D) & Label-free transformer: higher displacement error than the displacement-supervised one, lower errors on the other four metrics, each seed and the seed mean & 20 of 20 & 20 of 20 & replicated \\
H3 (2D) & A reading: energy-norm-supervised against label-free transformer, relative energy gap & $-$24.2\% (guard 38.1\%) & +99.1\% (guard 160.2\%) & no difference shown \\
\midrule
V & Per-instance values alike on the holdout and validation instances, 35 quantities & -- & none rejected & \hspace{0pt}representative \\
\bottomrule
\end{tabular}}

\smallskip
\begin{minipage}{\linewidth}\scriptsize\raggedright
Validation: the values of \cref{sec:norm,sec:labelfree,sec:checks} on the validation instances (T3 per load case). Holdout: the same stored networks on 686 (3D) and 4,111 (2D) held-out instances. T1 and D2: the comparisons that hold, five metrics in three seeds and the seed mean; T2: the lowest and highest share over the seeds, and the share of the instance-seed pairs; H1 to H3: rel and the noise guard of the note to \cref{tab:claims}; H2c: the rel of H2b, against $-$50\%. V: per network and metric, a two-sided Mann--Whitney test of the holdout against the validation instances, the 35 p-values taken together by Holm's procedure at a family-wise level of 5\% (supplementary Section~\SuppPrereg).
\end{minipage}
\end{table}


\Cref{tab:holdout} gives the claims with their values on the validation and on the held-out
instances; seven of the eight replicated. The eight verdicts are not independent: the claims
of each dimension read the same instances, and several of them the same networks; H2b and
H2c read one comparison, and the sign part of T3 is an identity (\cref{cor:zero}). As the
pre-registration states, agreement among dependent verdicts adds little evidence. In three
dimensions the label-free transformer was again more accurate than the supervised
transformer on all five metrics, in each seed and in the seed mean (T1); the graph network
again had the lower displacement error on most instances, \numHoldTTwoMin{} to
\numHoldTTwoMax{} of them in the three seeds, and the higher von Mises error on every
instance-seed pair (T2); and the supervised transformer and the graph network had load cases
worse than the zero field, \numHoldTThreeLab{} and \numHoldTThreeMgn{} of
\numHoldTThreeLoads{}, each with a positive energy and no other (T3). In two dimensions the
label-free transformer's lower relative energy gap (H1) and its larger displacement error
with smaller errors on the other four metrics (D2) replicated, as did the stiffness-norm
transformer's lower von Mises error, beyond the guard and by more than half (H2b and H2c).

The stiffness-norm transformer's lower relative energy gap (H2a) did not replicate. Its
seeds' relative energy gaps on the held-out instances were \numHoldHTwoANewSeeds{}, against
\numHoldHTwoABaseSeeds{} for the supervised transformer: \numHoldHTwoARel{} lower, within
the guard of \numHoldHTwoATau{}. The readings beside it, fixed in advance without a
criterion, locate the failure in a few instances. The stiffness-norm transformer's mean
relative energy gap was \numHoldHTwoARatio{} times as large on the held-out instances as on
the validation instances (95\% interval, the held-out instances resampled:
\numHoldHTwoARatioLow{} to \numHoldHTwoARatioHigh{}), whereas its median was much the same
on both, \numHoldHTwoAMedHold{} and \numHoldHTwoAMedVal{}. Of its \numHoldTwoPairs{}
instance-seed predictions on the held-out instances, \numHoldTwoWorseInstKnorm{} had a mean
relative energy gap above one, against \numHoldTwoWorseInstLab{} of the supervised
transformer's and \numHoldTwoWorseInstFree{} of the label-free transformer's, where on the
validation instances none of the three had any; of the \numHoldTwoLoads{} load cases,
\numHoldTwoWorseLoadKnorm{}, \numHoldTwoWorseLoadLab{} and \numHoldTwoWorseLoadFree{} were
worse than the zero field, each with a positive energy and no other. Read after the
evaluation, the largest relative energy gaps of the stiffness-norm transformer's three seeds
on the held-out instances were \numHoldHTwoATopSeeds{}. On the seeds' medians the reduction
against the supervised transformer was \numHoldHTwoAMedRel{}, beyond a guard of
\numHoldHTwoAMedTau{}. The claim of \cref{sec:results2d:norm} thus holds for the typical
instance, not for the mean, which a few held-out instances carry. H3, the stiffness-norm
transformer against the label-free transformer, again showed no difference beyond the guard;
on the seed means the label-free transformer's relative energy gap was now the lower,
\numHoldTwoFreeGap{} against \numHoldTwoKnormGap{}.

Beside the claims, the per-instance values themselves were compared, network by network and
metric by metric, between the held-out and the validation instances (V): of the
\numHoldVQuantities{} quantities none was rejected by a rank test at a family-wise level of
\numHoldVAlpha{}, and the validation instances were read as representative. That verdict
does not show that the means are unchanged: a mean that a few instances carry can move
without a rejection, as the stiffness-norm transformer's relative energy gap did.
\Cref{tab:holdoutvalues} gives the means and medians on both sets for every network and
metric, and \ref{app:prereg} the evaluation's two departures from its plan. The Ritz
projection of \cref{sec:methods:ritz}, added after the original analyses, was paired with the
networks on the held-out instances as on the validation instances: in three dimensions its
displacement error was smaller than the label-free transformer's on
\numRitzHoldFreeDispSixMin{} to \numRitzHoldFreeDispSixMax{} of the instances at $p=6$ and on
\numRitzHoldFreeDispNineMin{} to \numRitzHoldFreeDispNineMax{} at $p=9$, and its peak stress
error larger on at least \numRitzHoldFreePeakMin{} of them at every degree (supplementary
material, Table S1).
